\chapter{How to read a failed Test Run}\label{ch:failed-run}

A run that went wrong leaves four traces. Read them in this order, because the order is what makes the
cause visible.

\section{The run row first}

\texttt{test\_run.result} is \texttt{PASS} or \texttt{FAIL}. \texttt{end\_reason} is the story:

\begingroup\small
\begin{longtable}{p{3.2cm} p{9.4cm}}
\toprule
\textbf{End Reason} & \textbf{What happened} \\
\midrule
\endhead
\texttt{COMPLETED} & the schedule ran to its end, abort or not \\
\texttt{PRETEST\_FAILED} & a Pre-test Action raised; no step ran, and every step is NOTRUN \\
\texttt{STOPPED\_ON\_FAIL} & the Operator answered \texttt{ABORT} at a Stop-on-Fail prompt \\
\texttt{ABORTED} & a run-level \texttt{request\_abort()}, wherever it landed --- at a boundary, inside a
  step, or at a pause (packet 32b: the same button now gives the same reason) \\
\texttt{INTERRUPTED} & the worker ended on an exception it could not handle, or the process died with the
  run open and launch closed it \\
\bottomrule
\end{longtable}
\endgroup

A step's \texttt{result} is one of \texttt{PASS}, \texttt{FAIL}, \texttt{ERROR}, \texttt{SKIPPED},
\texttt{NOTRUN}, \texttt{FORCED\_NOTRUN} or \texttt{NOREPORT}. \texttt{NOREPORT} means the step ran and never
called \texttt{ATE.report} --- a schedule bug, recorded as a WARNING in the Run log rather than silently
counted as a pass.

\section{The two logs}

The \textbf{Run log} is the \texttt{log\_line} table: one row per line, with a level
(\texttt{DEBUG}/\texttt{INFO}/\texttt{WARNING}/\texttt{ERROR}), a source (\texttt{schedule}, \texttt{ate} or
\texttt{device}) and an owner --- the Action or the Attempt that produced it. Run log lines come from
\texttt{ATE.log}, redirected \texttt{print}, Device failures and warnings, and trace lines when tracing is
on. The \textbf{application log} is the per-launch file at \texttt{<data>/\allowbreak{}logs/\allowbreak{}ate-app-<time>.log}, which
is where an unhandled exception's traceback goes. They are different things and answer different questions:
``what did this run do'' versus ``what did the application do wrong''.

The last line of a step's Attempt in the report is the last \texttt{ATE.report} call of that Attempt; a
step whose Attempt raised records the error text as its \texttt{msg}.

\section{When the failure is the test harness, not the run}

If a test fails with ``a RunEngine worker outlived its test'' or ``N RunEngine worker(s) did not stop'',
the report section added by \texttt{describe\_\allowbreak{}stuck\_\allowbreak{}runs} is the thing to read: the live workers, what each
session thinks its state is, what the store has recorded, the store's write timings with any write still in
flight, and every thread's stack. Long or in-flight writes mean the disk. Normal writes with the worker's
stack elsewhere mean a fault the suite has not found yet. That is the diagnostic packet 32b added because
the Windows stall it could not reproduce had no explanation.

Read the warnings that follow an abort carefully: ``left open at end of run'' means a Device the schedule
never closed; ``was still closing when the aborted run ended'' means the close outlived the run's 3 s bound
and will finish later; ``datagram truncated to N byte(s)'' and ``timed out after N byte(s)'' are different
read failures, and the Device says which one it is.

\section{Reading the record with a query}

The fastest way to see what a run recorded, from a Station's database:

\begin{lstlisting}[language=ATEShell]
python -c "import sqlite3, sys; rows=sqlite3.connect(sys.argv[1]).execute('SELECT step_number,result,attempt_count FROM node ORDER BY node_seq'); print(rows.fetchall())" /tmp/station/data/ate.db
\end{lstlisting}

Read-only, no ATE involved. The database is passed as an argument rather than written inside the
string: on Windows the path has backslashes, and a path pasted into a string literal in
\texttt{-c} code is read as an escape sequence --- \texttt{C:\textbackslash Users\textbackslash ...}
is a \texttt{SyntaxError} before the query ever runs.

The same query against \texttt{log\_line} ordered by \texttt{seq} is the Run log
in emission order, which is the order the report shows.

\cando{go from a failed run to the row, the log line and the thread that produced it.}
